% =========================================================
% 6.4.2 DICCIONARIO DE DATOS
% Complementa los diagramas TikZ de 6.4.1 (er_*.tex) con lo que no
% muestran: trazabilidad, campos de significado no evidente, integridad
% referencial, restricciones y reglas del servidor. Los tipos, llaves y
% valores por defecto de cada campo están en los diagramas y no se repiten.
% Lo incluye seccion6_4.tex.
% =========================================================
% D{ancho}: columna p alineada a la izquierda (sin huecos de justificado)
\newcolumntype{D}[1]{>{\raggedright\arraybackslash}p{#1}}

\subsection{Diccionario de datos}
\label{sec:diccionario_datos}

Los tipos, las llaves y los valores por defecto de cada campo están en los diagramas de la sección anterior. Este apartado documenta lo que no muestran: el requerimiento que origina cada tabla, el significado de los campos que no se deduce de su nombre, las reglas de integridad referencial y las restricciones que se crean en la base de datos. Las referencias RN y RF remiten a las secciones \ref{sec:reglas_negocio} y \ref{sec:requerimientos}.

\subsubsection{Trazabilidad de las tablas}

% longtable: con su título no cabe en una página y una tabla [H] dejaría una página en blanco.
\begin{longtable}{|>{\footnotesize}c|>{\footnotesize}l|>{\footnotesize}D{0.35\textwidth}|>{\footnotesize}D{0.2\textwidth}|}
\caption{Propósito y origen de las tablas de la base de datos}\label{tab:dd_resumen} \\
\hline
\textbf{Nº} & \textbf{Tabla} & \textbf{Propósito} & \textbf{Origen} \\ \hline
\endfirsthead
\multicolumn{4}{c}{\bfseries Continuación de la Tabla \thetable} \\
\hline
\textbf{Nº} & \textbf{Tabla} & \textbf{Propósito} & \textbf{Origen} \\ \hline
\endhead
1 & \texttt{auth.\allowbreak users} & Cuentas de acceso, credenciales y verificación del correo (Supabase Auth) & RF-01, RF-02, RF-06 \\ \hline
2 & \texttt{auth.\allowbreak identities} & Formas de acceso vinculadas a una cuenta: correo o Google & RF-07 \\ \hline
3 & \texttt{estado\_cuenta} & Catálogo de estados de la cuenta & RF-12 \\ \hline
4 & \texttt{usuario} & Perfil y reputación de cada persona & RF-14, RN-10 \\ \hline
5 & \texttt{dispositivo} & Token de notificaciones y alertas por equipo & RF-42, RN-13 \\ \hline
6 & \texttt{solic\_eliminacion} & Plazo y constancia de la eliminación de la cuenta & RF-12, RN-23 \\ \hline
7 & \texttt{aviso\_privacidad} & Versiones del aviso de privacidad & RF-13, RN-24 \\ \hline
8 & \texttt{aviso\_privacidad\_usuario} & Versión del aviso que aceptó cada usuario & RF-13, RN-24 \\ \hline
9 & \texttt{manual} & Versiones del manual de usuario & RF-48 \\ \hline
10 & \texttt{categoria\_lugar} & Catálogo de categorías de lugar & RF-18 \\ \hline
11 & \texttt{favorito} & Lugares guardados por el usuario & RF-18 \\ \hline
12 & \texttt{ubicacion} & Última ubicación conocida del usuario, sin historial & RN-13, RN-24 \\ \hline
13 & \texttt{enlace} & Enlace y código que genera quien comparte su ubicación & RF-43, RF-44, RN-16 a RN-19 \\ \hline
14 & \texttt{estado\_seguimiento} & Catálogo de estados de un seguimiento & RF-45, RN-18 \\ \hline
15 & \texttt{seguimiento} & Personas que siguen a un usuario por medio de un enlace o un código & RF-45, RF-46, RN-18, RN-19 \\ \hline
16 & \texttt{categoria\_reporte} & Categorías de reporte y vigencia de cada estado & RN-07, RN-09 \\ \hline
17 & \texttt{reporte} & Reportes ciudadanos y su ciclo de vida & RF-37 a RF-39, RN-09 a RN-11 \\ \hline
18 & \texttt{valoracion} & Confirmaciones y desmentidos de cada reporte, por ciclo & RF-41, RN-11, RN-12 \\ \hline
19 & \texttt{navegacion} & Navegaciones del usuario y su evaluación & RF-29 a RF-35 \\ \hline
20 & \texttt{desglose} & Aristas que recorre cada navegación, en orden & RF-24, RF-25 \\ \hline
21 & \texttt{evento\_nav} & Desvíos, recálculos y alertas de pánico de una navegación & RF-32, RF-47, RN-20 \\ \hline
22 & \texttt{punto} & Vértices del grafo de la red peatonal & RF-22 \\ \hline
23 & \texttt{via} & Vías de OpenStreetMap de las que provienen las aristas & RF-22 \\ \hline
24 & \texttt{calle} & Aristas del grafo, con su trazo, su longitud y su riesgo & RF-22, RN-03 \\ \hline
25 & \texttt{emergencia} & Número de atención de cada corporación por zona & RF-47 \\ \hline
\end{longtable}

\subsubsection{Campos con significado no evidente}

\begin{table}[H]
\centering
\footnotesize
\caption{Campos cuyo significado no se deduce de su nombre}
\label{tab:dd_campos}
\renewcommand{\arraystretch}{1.3}
\begin{tabular}{|D{0.36\textwidth}|D{0.56\textwidth}|}
\hline
\textbf{Campo} & \textbf{Significado} \\ \hline
\texttt{usuario.\allowbreak reputacion} & Reputación de 0 a 5 estrellas. Determina cuántas confirmaciones necesitan los reportes del usuario (RN-10, RN-11) \\ \hline
\texttt{solic\_eliminacion.\allowbreak limite\_borrado} & Fecha máxima para borrar los datos personales, a siete días de la solicitud (RN-23) \\ \hline
\texttt{aviso\_privacidad.\allowbreak activo}, \texttt{manual.\allowbreak activo} & Marca la versión vigente. Solo una fila de cada tabla puede tenerlo en TRUE \\ \hline
\texttt{enlace.\allowbreak codigo}, \texttt{enlace.\allowbreak token\_acceso} & Las dos formas de compartir: el código de seis caracteres requiere que el seguido confirme (RN-18); el token viaja en el enlace directo de un solo uso (RN-17) \\ \hline
\texttt{enlace.\allowbreak activo}, \texttt{enlace.\allowbreak visible} & \texttt{activo} en FALSE pausa el enlace; \texttt{visible} en FALSE lo elimina de la lista del usuario sin borrar la fila \\ \hline
\texttt{enlace.\allowbreak id\_ubicacion} & Apunta a \texttt{ubicacion.\allowbreak id\_usuario}: para compartir su ubicación, el usuario necesita tener su fila en \texttt{ubicacion} \\ \hline
\texttt{seguimiento.\allowbreak ultima\_vista} & Última vez que el seguidor consultó la ubicación; permite avisar al seguido (RN-21) \\ \hline
\texttt{seguimiento.\allowbreak fecha\_confirmacion} & Momento en que el seguido aprobó la solicitud. Es NULL cuando se ingresó con un enlace \\ \hline
\texttt{reporte.\allowbreak estado} & Los valores RECHAZADO y EXPIRADO corresponden a los estados descartado y caducado de RN-09 y RN-11 \\ \hline
\texttt{reporte.\allowbreak ciclo}, \texttt{valoracion.\allowbreak ciclo} & Ronda de verificación. Aumenta cuando un reporte validado regresa a pendiente; los votos solo cuentan en el ciclo en que se emitieron (RN-09, RN-11) \\ \hline
\texttt{reporte.\allowbreak confirmaciones}, \texttt{reporte.\allowbreak rechazos} & Votos acumulados en el ciclo actual. Con tres rechazos el reporte pasa a RECHAZADO (RN-11) \\ \hline
\texttt{reporte.\allowbreak validaciones\_necesarias} & Confirmaciones necesarias para validar el reporte, según la reputación del autor (RN-11) \\ \hline
\texttt{reporte.\allowbreak fecha\_expiracion}, \texttt{reporte.\allowbreak expiracion\_aplicada} & Fin del periodo del estado actual, según la categoría (RN-09); la bandera indica que el proceso programado ya lo aplicó \\ \hline
\texttt{reporte.\allowbreak reputacion\_aplicada} & Evita modificar dos veces la reputación del autor por el mismo reporte (RN-10) \\ \hline
\texttt{valoracion.\allowbreak voto\_tipo} & 1 para una confirmación y $-1$ para un desmentido \\ \hline
\texttt{navegacion.\allowbreak franja\_horaria} & Franja con la que se calculó la ruta: diurna o nocturna (RN-04) \\ \hline
\texttt{desglose.\allowbreak secuencia} & Posición de la arista dentro de la ruta \\ \hline
\end{tabular}
\end{table}

\subsubsection{Integridad referencial}

Al eliminar una cuenta en \texttt{auth.\allowbreak users}, las llaves foráneas borran en cascada (\texttt{ON DELETE CASCADE}) el perfil y todos sus datos personales: dispositivos, aceptaciones del aviso, lugares favoritos, ubicación, enlaces con sus seguimientos, seguimientos como seguidor y navegaciones con sus aristas y eventos. Las excepciones son \texttt{reporte}, \texttt{valoracion} y \texttt{solic\_eliminacion}, que usan \texttt{ON DELETE SET NULL} para conservarse anonimizados (RN-23). Las llaves hacia los catálogos y hacia \texttt{aviso\_privacidad} usan \texttt{RESTRICT}, igual que \texttt{desglose.\allowbreak id\_calle}: no se borra un registro en uso. Las llaves de \texttt{calle} hacia \texttt{via} y \texttt{punto} no definen una acción, por lo que tampoco permiten borrar una vía o un vértice con aristas.

No existe una llave foránea entre \texttt{seguimiento} y \texttt{navegacion}. El seguidor llega a la ubicación y a la navegación del seguido por medio de \texttt{enlace.\allowbreak id\_ubicacion}, y una función de seguridad controla ese acceso.

\subsubsection{Restricciones e índices}

Las restricciones de la tabla \ref{tab:dd_check} no se ven en el diagrama y se crean en la base de datos. Además, se crean índices GiST sobre \texttt{reporte.\allowbreak ubicacion}, \texttt{ubicacion.\allowbreak ubicacion} y \texttt{emergencia.\allowbreak poligono} para las búsquedas por cercanía, y sobre los extremos de las aristas (\texttt{calle.\allowbreak id\_origen} y \texttt{calle.\allowbreak id\_destino}) para recorrer el grafo.

\begin{table}[H]
\centering
\footnotesize
\caption{Restricciones CHECK e índices únicos parciales}
\label{tab:dd_check}
\renewcommand{\arraystretch}{1.3}
\begin{tabular}{|D{0.25\textwidth}|D{0.66\textwidth}|}
\hline
\textbf{Tabla} & \textbf{Regla} \\ \hline
\texttt{usuario} & \texttt{reputacion} entre 0 y 5 (RN-10) \\ \hline
\texttt{solic\_eliminacion} & \texttt{estado} en (PENDIENTE, EJECUTADA); \texttt{limite\_borrado} a no más de siete días de \texttt{fecha\_solic} (RN-23); \texttt{ejecutada\_en} con valor solo si el estado es EJECUTADA \\ \hline
\texttt{aviso\_privacidad}, \texttt{manual} & Índice único parcial: una sola fila con \texttt{activo} en TRUE \\ \hline
\texttt{enlace} & \texttt{vigencia} mayor que \texttt{fecha\_alta} \\ \hline
\texttt{categoria\_reporte} & \texttt{duracion\_estado} mayor que cero \\ \hline
\texttt{reporte} & \texttt{confirmaciones} y \texttt{rechazos} mayores o iguales a cero; \texttt{validaciones\_necesarias} entre 1 y 3 (RN-11); \texttt{fecha\_expiracion} mayor que \texttt{fecha\_alta} \\ \hline
\texttt{valoracion} & \texttt{voto\_tipo} en (1, $-1$) \\ \hline
\texttt{navegacion} & \texttt{calificacion} entre 1 y 5; \texttt{franja\_horaria} en (DIURNA, NOCTURNA) (RN-04). Índice único parcial: una sola navegación EN\_CURSO por usuario \\ \hline
\texttt{calle} & \texttt{riesgo} entre 0 y 1; \texttt{distancia\_metros} mayor que cero \\ \hline
\end{tabular}
\end{table}

\subsubsection{Reglas que aplica el servidor}

Las reglas de la tabla \ref{tab:dd_reglas_servidor} se implementan con funciones, procesos programados y políticas de seguridad (RLS) de la base de datos, no con su estructura.

\begin{table}[H]
\centering
\footnotesize
\caption{Reglas de negocio que aplica el servidor}
\label{tab:dd_reglas_servidor}
\renewcommand{\arraystretch}{1.3}
\begin{tabular}{|D{0.5\textwidth}|D{0.41\textwidth}|}
\hline
\textbf{Regla} & \textbf{Mecanismo} \\ \hline
Ubicación del reporte dentro del área de servicio (RN-01) & Función de registro de reportes \\ \hline
Confirmaciones necesarias según la reputación del autor (RN-11) & Función que asigna \texttt{validaciones\_necesarias} al iniciar cada ciclo \\ \hline
Voto de usuarios distintos del autor, a 1 km o menos, uno por ciclo (RN-11, RN-12) & Función de votación y restricción única de \texttt{valoracion} \\ \hline
Validación y rechazo del reporte (RN-11) & Función de evaluación de votos \\ \hline
Fin del periodo de cada estado y regreso de validado a pendiente (RN-09) & Proceso programado sobre \texttt{fecha\_expiracion} \\ \hline
Reputación modificada una sola vez por reporte (RN-10) & Función que consulta \texttt{reputacion\_aplicada} \\ \hline
Avisos de reportes cercanos (RN-13) & Webhook y función que consulta \texttt{ubicacion} y \texttt{dispositivo} \\ \hline
Acceso del seguidor y aviso al seguido (RN-17 a RN-19, RN-21) & Función de seguridad que exige un seguimiento APROBADO y un enlace activo, visible y vigente, y actualiza \texttt{ultima\_vista} \\ \hline
Borrado de los datos personales en el plazo de la solicitud (RN-23) & Proceso programado sobre \texttt{limite\_borrado} \\ \hline
Catálogos de solo lectura; votos y seguimientos sin inserción directa & Políticas RLS: los usuarios solo leen los catálogos y solo las funciones insertan votos y seguimientos \\ \hline
\end{tabular}
\end{table}
